\clearpage\phantomsection\label{exercise-index-3}
\BookHeading{Percorso degli esercizi guidati}
Quindici esercizi accompagnano lo studio del modulo. Ogni traccia segue le sezioni richieste e rinvia a uno svolgimento completo. Si suggerisce di formulare la previsione e tentare i passaggi prima di consultare la soluzione.
\begin{longtable}{@{}p{13mm}p{75mm}r@{}}\toprule Codice & Argomento & Pagina\\\midrule\endhead
\hyperref[ex-G3.01]{G3.01} & Tradurre un circuito nella sequenza delle operazioni & \pageref{ex-G3.01}\\[3pt]
\hyperref[ex-G3.02]{G3.02} & Costruire la versione reversibile di una piccola funzione & \pageref{ex-G3.02}\\[3pt]
\hyperref[ex-G3.03]{G3.03} & Applicare \(X\), \(Z\) e \(H\) agli stessi stati & \pageref{ex-G3.03}\\[3pt]
\hyperref[ex-G3.04]{G3.04} & Preparare uno stato assegnato & \pageref{ex-G3.04}\\[3pt]
\hyperref[ex-G3.05]{G3.05} & Costruire e verificare l'inversa di un circuito & \pageref{ex-G3.05}\\[3pt]
\hyperref[ex-G3.06]{G3.06} & Calcolare un interferometro con fase \(\pi /3\) & \pageref{ex-G3.06}\\[3pt]
\hyperref[ex-G3.07]{G3.07} & Applicare una porta al filo corretto & \pageref{ex-G3.07}\\[3pt]
\hyperref[ex-G3.08]{G3.08} & Seguire CNOT su base e sovrapposizione & \pageref{ex-G3.08}\\[3pt]
\hyperref[ex-G3.09]{G3.09} & Preparare e misurare uno stato di Bell & \pageref{ex-G3.09}\\[3pt]
\hyperref[ex-G3.10]{G3.10} & Usare e ripulire un'ancilla & \pageref{ex-G3.10}\\[3pt]
\hyperref[ex-G3.11]{G3.11} & Verificare una porta composta & \pageref{ex-G3.11}\\[3pt]
\hyperref[ex-G3.12]{G3.12} & Svolgere i quattro casi di Deutsch & \pageref{ex-G3.12}\\[3pt]
\hyperref[ex-G3.13]{G3.13} & Eseguire una ricerca fra quattro alternative & \pageref{ex-G3.13}\\[3pt]
\hyperref[ex-G3.14]{G3.14} & Seguire un'evoluzione e la sua inversa & \pageref{ex-G3.14}\\[3pt]
\hyperref[ex-G3.15]{G3.15} & Verificare un programma completo & \pageref{ex-G3.15}\\[3pt]
\bottomrule\end{longtable}
{\small I primi cinque esercizi offrono una guida esplicita; i successivi cinque consolidano il calcolo; gli ultimi cinque collegano più concetti. Tutti conservano uno svolgimento completo.}\par
